\section{Всестороннее оисследование по безопасному обучению с подкреплением}

\subsection{Ключевые идеи}

В статье~\cite{garcia2015} предложена систематизация и формальная классификация методов безопасного обучения с подкреплением.

Авторы выделяют две фундаментальные категории подходов:
\begin{itemize}
    \item \emph{Модификация критерия оптимальности}~-- добавление меры риска или фактора безопасности в классическую функцию оптимизации.
    \item \emph{Модификация процесса исследования среды}~-- изменение стратегии исследования посредством включения внешних знаний, советов «учителя» или метрики риска.
\end{itemize}

Цель работы~-- создание единой таксономии и аналитической основы для систематического сравнения существующих методов в безопасном обучении с подкреплением, а также формулирование направлений дальнейших исследований.
\subsection{Математическое описание метода}

Авторы выделяют четыре типа критериев, которые по-разному определяют понятие безопасности в обучении с подкреплением.  
Эти подходы составляют первую категорию классификации~-- методы, \emph{модифицирующие критерий оптимальности}.

\begin{itemize}

    \item \emph{Риск-нейтральный критерий.}  
    Агент максимизирует ожидаемое дисконтированное вознаграждение:
    \begin{equation}
    \max_{\pi \in \Pi} \mathbb{E}_{\pi}\left[\sum_{t=0}^{\infty} \gamma^t r_t \right],
    \end{equation}
    без учёта возможных редких, но катастрофических событий.

    \item \emph{Критерий наихудшего случая.}  
    Оптимизация проводится с учётом наихудшего сценария среды:
    \begin{equation}
    \max_{\pi \in \Pi} \min_{w \in \Omega_\pi} \mathbb{E}_{\pi, w}\left[\sum_{t=0}^{\infty} \gamma^t r_t \right],
    \end{equation}
    что обеспечивает робастность политики к неопределённости.

    \item \emph{Риск-чувствительный критерий.}  
    Функция полезности учитывает отношение агента к риску:
    \begin{equation}
    \max_{\pi \in \Pi} \frac{1}{\beta} \log \mathbb{E}_{\pi}\left[\exp\left(\beta \sum_{t=0}^{\infty} \gamma^t r_t\right)\right],
    \end{equation}
    где параметр $\beta$ определяет степень избегания или принятия риска.

    \item \emph{Критерий с ограничениями.}  
    Агент максимизирует ожидаемое вознаграждение при соблюдении условий безопасности:
    \begin{equation}
    \begin{aligned}
    \max_{\pi \in \Pi} & \quad \mathbb{E}_{\pi}[R], \\
    \text{при этом} & \quad c_i(\pi) \leq \alpha_i, \quad i = 1,\dots,n.
    \end{aligned}
    \end{equation}

\end{itemize}

\subsection{Алгоритм}

Работа имеет обзорный характер, поэтому авторы не предлагают нового алгоритма, а систематизируют существующие подходы безопасного обучения с подкреплением.  
Они выделяют две основные категории методов, отражающие разные способы достижения безопасности.

\begin{itemize}

    \item \emph{Методы, модифицирующие критерий оптимальности.}  
    Безопасность достигается через изменение целевой функции агента, включающей риск или неопределённость среды.  
    Сюда относятся:
    \begin{itemize}
        \item методы наихудшего случая~-- обеспечивают робастность политики, выбирая стратегию, устойчивую к неблагоприятным сценариям;
        \item риск-чувствительные методы~-- вводят меру риска в функцию полезности, регулируя баланс между вознаграждением и стабильностью;
        \item методы с ограничениями~-- формализуют безопасность через ограничения на вероятность ошибок или уровень ущерба.
    \end{itemize}

    \item \emph{Методы, модифицирующие процесс исследования среды.}  
    Безопасность обеспечивается управлением исследованием, чтобы агент избегал опасных состояний во время обучения.  
    В таких подходах агент использует внешние знания, демонстрации безопасного поведения или корректирующие сигналы, ограничивающие риск.  
    Эти методы направлены на предотвращение катастрофических ситуаций и повышение надёжности обучения в реальных условиях.

\end{itemize}


\subsection{Наборы данных}

Поскольку большинство алгоритмов безопасного обучения с подкреплением проверяются в симулированных условиях, авторы рассматривают несколько типовых сред, моделирующих риск и возможность катастрофических состояний.  
Они используются для оценки устойчивости и надёжности стратегий агента при обучении.

\begin{itemize}

    \item \emph{Cliff World}~-- двумерная дискретная среда с областями, ведущими к катастрофическим потерям вознаграждения.  
    Применяется для проверки способности алгоритмов избегать опасных состояний и балансировать между скоростью и безопасностью~\cite{gaskett2003}.

    \item \emph{Grid World с ограничениями}~-- сеточная среда с зонами повышенного риска, где часть состояний или действий недоступны.  
    Используется для тестирования методов, интегрирующих внешние знания и ограничения в процесс обучения~\cite{driessens2004}.

    \item \emph{Pole Balancing (Cart-Pole)}~-- непрерывная задача стабилизации, где ошибка управления приводит к падению и большим отрицательным вознаграждениям.  
    Применяется для анализа устойчивости и риск-чувствительности алгоритмов~\cite{geibel2005}.

    \item \emph{Helicopter Control Task}~-- сложная динамическая среда управления вертолётом, где безопасность имеет физический смысл.  
    Используется для оценки методов, обеспечивающих безопасное обучение и имитацию поведения~\cite{abbeel2005, abbeel2010}.

    \item \emph{Mountain Car} и \emph{Acrobot}~-- классические тестовые задачи, моделирующие физические ограничения и потенциальные «опасные» состояния.  
    Применяются для проверки алгоритмов, учитывающих риск и ограничения на действия~\cite{mihatsch2002, moldovan2012}.

\end{itemize}


\subsection{Метрики}

Для оценки эффективности и безопасности алгоритмов обучения с подкреплением используются несколько метрик, отражающих как успешность политики, так и уровень риска при обучении.  
Авторы подчёркивают, что надёжная оценка достигается только при совместном анализе нескольких показателей.

\begin{itemize}

    \item \emph{Ожидаемое суммарное вознаграждение}~$\mathbb{E}[R]$.  
    Основная метрика, отражающая среднее качество политики:
    \begin{equation}
    \mathbb{E}[R] = \mathbb{E}_\pi \left[\sum_{t=0}^{\infty} \gamma^t r_t \right],
    \end{equation}
    где $\gamma$~-- коэффициент дисконтирования.  
    Используется для оценки эффективности стратегии при сохранении ограничений безопасности~\cite{mihatsch2002, geibel2005}.

    \item \emph{Дисперсия вознаграждения}~$\mathrm{Var}(R)$.  
    Характеризует стабильность поведения агента: чем меньше дисперсия, тем устойчивее результаты.  
    В риск-чувствительных методах она включается в целевую функцию как мера избегания риска~\cite{mihatsch2002, borkar2002}.

    \item \emph{Вероятность попадания в ошибочные состояния}~$P(s \in S_{err})$.  
    Отражает частоту нарушений ограничений безопасности, например столкновений или выходов за допустимые границы.  
    Часто используется как ограничение в оптимизационной задаче:
    \begin{equation}
    \mathbb{E}_\pi[c(s,a)] \leq \alpha,
    \end{equation}
    где $\alpha$~-- максимально допустимая вероятность ошибок~\cite{geibel2005, moldovan2012}.

    \item \emph{Частота катастрофических событий}.  
    Показывает, насколько часто агент совершает критические ошибки, приводящие к значительным потерям.  
    Используется для оценки качества обучения в средах с высоким риском~\cite{gaskett2003, abbeel2005}.

    \item \emph{Надёжность (robustness score)}.  
    Оценивает устойчивость политики к изменению параметров среды и неопределённости переходов.  
    Рассчитывается как худшее ожидаемое вознаграждение по множеству моделей среды:
    \begin{equation}
    R_{robust} = \min_{T \in \mathcal{T}} \mathbb{E}_{\pi, T}[R],
    \end{equation}
    и используется для анализа робастных стратегий~\cite{heger1994, moldovan2012}.

\end{itemize}


\subsection{Модели}

Все подходы, рассматриваемые авторами, формулируются в рамках \emph{марковского процесса принятия решений}, задаваемого кортежем
\begin{equation}
\langle S, A, T, R \rangle,
\end{equation}
где $S$~-- множество состояний, $A$~-- множество действий, $T(s'|s,a)$~-- функция переходов, а $R(s,a)$~-- функция вознаграждения.  
Базовая цель агента формулируется как максимизация ожидаемого дисконтированного возврата:
\begin{equation}
\max_{\pi \in \Pi} \mathbb{E}_{\pi}\left[\sum_{t=0}^{\infty} \gamma^t r_t \right].
\end{equation}

Далее авторы рассматривают расширения этой модели, отражающие разные способы учёта безопасности и неопределённости.

\begin{itemize}

    \item \emph{Риск-нейтральная модель.}  
    Описывает стандартное обучение с подкреплением, где агент максимизирует среднее вознаграждение без учёта риска.  
    Эта модель используется как базовая для сравнения с безопасными и робастными вариантами.

    \item \emph{Робастная модель с наихудшим случаем.}  
    Предполагает неопределённость в переходах среды.  
    Оптимизация проводится по принципу наихудшего сценария:
    \begin{equation}
    \max_{\pi \in \Pi} \min_{T \in \mathcal{T}} \mathbb{E}_{\pi,T}\left[\sum_{t=0}^{\infty} \gamma^t r_t \right],
    \end{equation}
    что обеспечивает устойчивость политики к ошибкам моделирования и изменениям динамики среды.

    \item \emph{Риск-чувствительная модель.}  
    Использует экспоненциальную функцию полезности для учёта отношения к риску:
    \begin{equation}
    \max_{\pi \in \Pi} \frac{1}{\beta} \log \mathbb{E}_{\pi}\left[\exp\left(\beta \sum_{t=0}^{\infty} \gamma^t r_t\right)\right],
    \end{equation}
    где параметр $\beta$ определяет степень склонности или избегания риска.

    \item \emph{Ограниченная модель.}  
    Включает явные ограничения безопасности:
    \begin{equation}
    \begin{aligned}
    \max_{\pi \in \Pi} & \quad \mathbb{E}_{\pi}[R], \\
    \text{при этом} & \quad \mathbb{E}_{\pi}[c_i(s,a)] \leq \alpha_i.
    \end{aligned}
    \end{equation}
    Здесь функции $c_i$ описывают риск или нарушения безопасности, а $\alpha_i$~-- допустимые пороги.

\end{itemize}

Все методы безопасного обучения с подкреплением интерпретируются как расширения базовой модели марковского процесса принятия решений, в которых безопасность вводится через робастность, чувствительность к риску или ограничения на поведение агента.
